% !PRINTED 594
\weberpage{594}
%<<CONT
noch den \gesperrt{Fermat'schen} Lehrsatz für rationale Zahlen an, so erkennt man, dass die Congruenzen (2) und (5) \gesperrt{gleichbedeutend} sind mit den für \gesperrt{jede Zahl} $\omega$ in $\mathfrak{o}$ gültigen Formeln

\weq{6}{\omega\mkern1mu|\mkern1mu\chi \equiv \omega,\quad \omega^P \equiv \omega\mkern1mu|\mkern1mu\psi_0\ (\mathrm{mod}\ \mathfrak{p}).}
Aus der zweiten Congruenz (6) folgt, dass, wenn $\omega$ durch $\mathfrak{p}$ theilbar ist, immer auch $\omega\mkern1mu|\mkern1mu\psi_0$ durch $\mathfrak{p}$ theilbar sein muss. Folglich ist $\mathfrak{p}$ durch $\mathfrak{p}\mkern1mu|\mkern1mu\psi_0$ theilbar, und als Primideal $= \mathfrak{p}\mkern1mu|\mkern1mu\psi_0$. Daraus folgt, dass $\psi_0$ in der Gruppe $\Psi$ enthalten ist.

Wir verstehen jetzt unter $\gamma$ eine Primitivwurzel von $\mathfrak{p}$ und wenden die am Schlusse des \S.\,161 bewiesene Formel an:
% [sic: The source repeats $\gamma^{P^f-1}$ in the last factor, as on printed page 592; the P-power orbit instead requires $\gamma^{P^{f-1}}$. The printed expression is retained.]
\weq{7}{(t-\gamma)(t-\gamma^P)\ldots(t-\gamma^{P^f-1}) \equiv F(t)\ (\mathrm{mod}\ \mathfrak{p}),}
in der $F(t)$ eine ganze Function von $t$ in $R$ ist.

Daraus folgt
\[ F(\gamma) \equiv 0\ (\mathrm{mod}\ \mathfrak{p}), \]
und wenn nun $\psi$ irgend eine Substitution aus $\Psi$ ist:
\[ F(\gamma\mkern1mu|\mkern1mu\psi) \equiv 0\ (\mathrm{mod}\ \mathfrak{p}). \]

Daraus ergiebt sich aber, dass $\gamma\mkern1mu|\mkern1mu\psi$ mit einer der in (7) vorkommenden Potenzen von $\gamma$ nach $\mathfrak{p}$ congruent sein muss, also etwa
\weq{8}{\gamma^{P^\nu} \equiv \gamma\mkern1mu|\mkern1mu\psi\ (\mathrm{mod}\ \mathfrak{p}).}
Nun ist jede durch $\mathfrak{p}$ nicht theilbare Zahl $\omega$ in $\mathfrak{o}$ mit einer Potenz von $\gamma$ congruent, und wenn man also (8) zu dieser Potenz erhebt, so folgt
\weq{9}{\omega^{P^\nu} \equiv \omega\mkern1mu|\mkern1mu\psi\ (\mathrm{mod}\ \mathfrak{p}),}
und diese Congruenz gilt offenbar auch noch, wenn $\omega$ durch $\mathfrak{p}$ theilbar ist, also für \gesperrt{alle Zahlen} in $\mathfrak{o}$.

Wenn wir nun die zweite der Congruenzen (6) $\nu$ mal nach einander anwenden, so folgt
\weq{10}{\omega^{P^\nu} \equiv \omega\mkern1mu|\mkern1mu\psi_0^\nu\ (\mathrm{mod}\ \mathfrak{p}),}
also
\weq{11}{\omega\mkern1mu|\mkern1mu\psi \equiv \omega\mkern1mu|\mkern1mu\psi_0^\nu\ (\mathrm{mod}\ \mathfrak{p}),}
und wenn wir $\omega$ durch $\omega\mkern1mu|\mkern1mu\psi^{-1}$ ersetzen oder auf (11) die Substitution $\psi^{-1}$ anwenden:
\weq{12}{\omega \equiv \omega\mkern1mu|\mkern1mu\psi^{-1}\psi_0^\nu \equiv \omega\mkern1mu|\mkern1mu\psi_0^\nu\psi^{-1}\ (\mathrm{mod}\ \mathfrak{p}).}
Es sind also sowohl $\psi^{-1}\psi_0^\nu$ als auch $\psi_0^\nu\psi^{-1}$ in $X$ enthalten,
%CONT>>
